% !TEX root = ../matan.tex

\section{Локальные экстремумы. Определение и необходимое условие экстремума. Стационарные точки}

\subsection*{Определение локального экстремума}

\begin{Definition}[required]{Локальный экстремум}{}
	Пусть $f\colon\Omega\to\RR$, $\Omega\subset\RR^d$, и $x$ --- внутренняя точка
	области $\Omega$. Точка $x$ называется \textbf{точкой локального экстремума}
	функции $f$, если существует окрестность $U_x\subset\Omega$, в которой приращение
	$f(\xi)-f(x)$ сохраняет знак для всех $\xi\in U_x$:
	\begin{itemize}
		\item $f(\xi)\le f(x)$ --- \emph{локальный максимум} (строгий при $f(\xi)<f(x)$,
		$\xi\ne x$);
		\item $f(\xi)\ge f(x)$ --- \emph{локальный минимум} (строгий при $f(\xi)>f(x)$,
		$\xi\ne x$).
	\end{itemize}
\end{Definition}

\subsection*{Необходимое условие экстремума}

\begin{Theorem}{Необходимое условие локального экстремума}{}
	Пусть $x$ --- точка локального экстремума функции $f\colon\Omega\to\RR$, и в этой
	точке существует градиент $\nabla f(x)$ (то есть существуют все частные
	производные). Тогда
	\[
	\nabla f(x)=\bigl(\partial_1 f(x),\dots,\partial_d f(x)\bigr)=\vec 0.
	\]
\end{Theorem}

\begin{proof}
	Зафиксируем индекс $i\in\{1,\dots,d\}$ и все переменные, кроме $i$-й. Рассмотрим
	вспомогательную функцию одной переменной
	\[
	g_i(t)=f(x_1,\dots,x_{i-1},\,x_i+t,\,x_{i+1},\dots,x_d),
	\]
	определённую при достаточно малых $|t|$ (так что аргумент остаётся в $U_x$).
	Поскольку $x$ --- точка локального экстремума $f$, приращение
	$g_i(t)-g_i(0)=f(\dots,x_i+t,\dots)-f(x)$ сохраняет знак при малых $t$, то есть
	$t=0$ --- точка локального экстремума функции $g_i$.

	Существование частной производной $\partial_i f(x)$ означает дифференцируемость
	$g_i$ в нуле, причём $g_i'(0)=\partial_i f(x)$. По теореме Ферма (необходимое
	условие экстремума функции одной переменной) $g_i'(0)=0$, поэтому
	$\partial_i f(x)=0$. Повторяя рассуждение для всех $i=1,\dots,d$, получаем
	$\nabla f(x)=\vec 0$.
\end{proof}

\subsection*{Стационарные точки}

\begin{Definition}{Стационарная и регулярная точки}{}
	\begin{itemize}
		\item Точка $x$, в которой $\nabla f(x)=\vec 0$, называется
		\textbf{стационарной точкой} функции $f$.
		\item Стационарная точка называется \textbf{регулярной (невырожденной)}, если
		$f$ дважды дифференцируема в точке $x$ и определитель её матрицы Гессе
		(матрицы вторых частных производных) отличен от нуля:
		\[
		\det\Bigl(\bigl\{\partial_k\partial_j f(x)\bigr\}_{k,j=1}^{d}\Bigr)\ne 0.
		\]
	\end{itemize}
\end{Definition}

В этих терминах необходимое условие звучит так: \emph{всякая точка локального
экстремума, в которой существует градиент, является стационарной.} Поиск экстремумов
поэтому начинают с решения системы $\nabla f(x)=\vec 0$.

\begin{Remark}{Условие не является достаточным}{}
	Обратное неверно: стационарная точка может не быть экстремумом. Для
	$f(x_1,x_2)=x_1^2-x_2^2$ имеем $\nabla f(0,0)=\vec 0$, однако $f(t,0)=t^2>0$ и
	$f(0,t)=-t^2<0$ в любой окрестности нуля --- это \emph{седловая точка}. Различить
	экстремум и седло позволяют достаточные условия (знакоопределённость второго
	дифференциала), которые применимы именно в регулярных стационарных точках.
\end{Remark}
